\booksubheading{File কেন প্রয়োজন?}

\begin{definitionbox}

\textbf{File:}
কম্পিউটারের Secondary Storage-এ একটি নির্দিষ্ট নামের অধীনে
Data স্থায়ীভাবে বা ভবিষ্যতে ব্যবহারের জন্য সংরক্ষণের একককে
\textbf{File} বলে।

\end{definitionbox}

Variable Program চলাকালীন Data সংরক্ষণ করতে পারে।

\begin{codeoutputbox}{Variable-এ Data}

\noindent
\begin{minipage}[t]{0.50\linewidth}
\codeheading

\begin{english}
\ttfamily
marks = 85\\
print(marks)
\end{english}

\end{minipage}
\hfill
\begin{minipage}[t]{0.46\linewidth}
\outputheading

\begin{center}
\textenglish{\ttfamily 85}
\end{center}

\textenglish{\ttfamily marks}
Variable Program চলার সময় Value ধরে রেখেছে।

কিন্তু Program বন্ধ হওয়ার পর Runtime Memory-এর এই Data
স্থায়ীভাবে থাকে না।

\end{minipage}

\end{codeoutputbox}

Program বন্ধ হওয়ার পরও Data প্রয়োজন হলে তা স্থায়ী Storage-এ
সংরক্ষণ করতে হয়। এ কাজে \textbf{File} ব্যবহার করা যায়।

\begin{rememberbox}

\begin{center}

\textbf{Program চলাকালীন Data}
\(\longrightarrow\)
\textbf{Variable}

\medskip

\textbf{Program বন্ধ হওয়ার পরও Data দরকার}
\(\longrightarrow\)
\textbf{File}

\end{center}

\end{rememberbox}

%================================================
% Temporary Persistent
%================================================

\booksubheading{Temporary ও Persistent Data}

\begin{definitionbox}

\textbf{Temporary Data:}
Program চলাকালীন সাময়িকভাবে Memory-তে থাকা এবং সাধারণত Program
শেষ হলে হারিয়ে যাওয়া Data-কে \textbf{Temporary Data} বলা হয়।

\medskip

\textbf{Persistent Data:}
Program শেষ হওয়ার পরও Storage-এ সংরক্ষিত থাকে এবং ভবিষ্যতে
পুনরায় ব্যবহার করা যায় এমন Data-কে \textbf{Persistent Data} বলা হয়।

\end{definitionbox}

\begin{tabularx}
{\linewidth}
{
|>{\bfseries\RaggedRight\arraybackslash}p{0.29\linewidth}
|Y
|Y|
}

\hline
\rowcolor{softBlue}

\textbf{বিষয়}
&
\textbf{Variable / Runtime Memory}
&
\textbf{File}
\\
\hline

Data-এর প্রকৃতি
&
Temporary
&
Persistent
\\
\hline

কতক্ষণ থাকে
&
Program চলাকালে
&
Program বন্ধ হওয়ার পরও থাকতে পারে
\\
\hline

প্রধান ব্যবহার
&
Processing
&
Data সংরক্ষণ
\\
\hline

পরে ব্যবহার
&
Runtime শেষ হলে থাকে না
&
পুনরায় ব্যবহার করা যায়
\\
\hline

উদাহরণ
&
Calculation Result
&
Student Record, Result
\\
\hline

\end{tabularx}

\begin{rememberbox}

\begin{center}

\textbf{Variable}
\(\longrightarrow\)
\textbf{Temporary Runtime Data}

\qquad

\textbf{File}
\(\longrightarrow\)
\textbf{Persistent Data}

\end{center}

\end{rememberbox}

%================================================
% File uses
%================================================

\booksubheading{File ব্যবহারের প্রয়োজনীয়তা}

File ব্যবহার করা হয়—

\begin{compacttwocolpoints}

\item Program বন্ধ হওয়ার পরও Data সংরক্ষণ করতে।

\item পরবর্তীতে Data আবার পড়তে ও ব্যবহার করতে।

\item Student, Customer বা অন্যান্য Record রাখতে।

\item Result, Report বা Log সংরক্ষণ করতে।

\item অনেক Data দীর্ঘ সময় ধরে রাখতে।

\item Saved Data ভবিষ্যতে আবার Process করতে।

\end{compacttwocolpoints}

\booksubheading{Python-এ File Handling}

Python Program শুধু Runtime-এর Variable নিয়ে কাজ করে না; প্রয়োজন হলে
Storage-এ থাকা File-এ Data লিখতে, File থেকে Data পড়তে এবং নতুন Data
যোগ করতেও পারে। Program-এর মাধ্যমে File খোলা, Data পড়া/লেখা এবং
File বন্ধ করার পুরো প্রক্রিয়াকে \textbf{File Handling} বলা হয়।

\begin{definitionbox}
\textbf{\color{darkTeal}File Handling:}

Python Program-এর সাহায্যে কোনো File \textbf{Open করা, Data Read করা,
Data Write/Append করা এবং কাজ শেষে File Close করার প্রক্রিয়াকে}
\textbf{File Handling} বলা হয়।
\end{definitionbox}

\begin{center}
\textbf{Python Program}
\(\longrightarrow\)
\textbf{\textenglish{\ttfamily open()}}
\(\longrightarrow\)
\textbf{File Object}
\(\longrightarrow\)
\textbf{Read / Write / Append}
\(\longrightarrow\)
\textbf{Close}
\end{center}

\begin{rememberbox}
\textbf{\color{deepOrange}Python-এর সঙ্গে File-এর সম্পর্ক}

Python হলো Program লেখার Language; আর File হলো Data স্থায়ীভাবে রাখার
Storage মাধ্যম। Python Program \textenglish{\ttfamily open()} Function-এর
মাধ্যমে File-এর সঙ্গে সংযোগ তৈরি করে এবং সেই File Object ব্যবহার করে
Data পড়া বা লেখা হয়।
\end{rememberbox}

%================================================
% Program file vs data file
%================================================

\booksubheading{Python Program File ও Data File এক নয়}

\begin{tabularx}{\linewidth}{|>{\bfseries\RaggedRight\arraybackslash}p{0.28\linewidth}|Y|Y|}
\hline
\rowcolor{softBlue}
\textbf{বিষয়} & \textbf{Python Program File} & \textbf{Data File}\\
\hline
মূল কাজ & Python Code সংরক্ষণ & Program-এর প্রয়োজনীয় Data সংরক্ষণ\\
\hline
সাধারণ Extension & \textenglish{\ttfamily .py} & \textenglish{\ttfamily .txt}, \textenglish{\ttfamily .csv} ইত্যাদি\\
\hline
উদাহরণ & \textenglish{\ttfamily result.py} & \textenglish{\ttfamily result.txt}\\
\hline
সম্পর্ক & Code Execute করে & Code-এর মাধ্যমে Read/Write করা যায়\\
\hline
\end{tabularx}

\begin{center}
\textenglish{\ttfamily result.py}
\(\longrightarrow\)
\textbf{Python Program}
\(\longrightarrow\)
\textenglish{\ttfamily result.txt}
\end{center}

অর্থাৎ \textenglish{\ttfamily .py} File-এ Program লেখা থাকে, আর
\textenglish{\ttfamily .txt}-এর মতো Data File-এ Program-এর Data
সংরক্ষণ করা যেতে পারে।

%================================================
% open
%================================================

\booksubheading{\textenglish{\ttfamily open()} দিয়ে File খোলা}

Python-এ File ব্যবহার করার আগে সাধারণত \textenglish{\ttfamily open()}
Function দিয়ে File খোলা হয়।

\begin{center}
\textenglish{\ttfamily file = open("filename", "mode")}
\end{center}

এখানে—

\begin{itemize}
\item \textenglish{\ttfamily "filename"} হলো যে File নিয়ে কাজ করা হবে তার নাম বা Path।
\item \textenglish{\ttfamily "mode"} দিয়ে File-টি Read, Write নাকি Append করার জন্য খোলা হবে তা নির্ধারণ করা হয়।
\item \textenglish{\ttfamily open()} একটি \textbf{File Object} Return করে; এই Object-এর মাধ্যমে File-এর Data নিয়ে কাজ করা হয়।
\end{itemize}

\begin{codeoutputbox}{File Open ও Close করার মৌলিক গঠন}
\noindent
\begin{minipage}[t]{0.54\linewidth}
\vspace{0pt}
\codeheading
\begin{english}
\ttfamily
file = open("result.txt", "w")\\
file.write("ICT = 85")\\
file.close()
\end{english}
\end{minipage}
\hfill
\begin{minipage}[t]{0.42\linewidth}
\vspace{0pt}
\outputheading

\textenglish{\ttfamily result.txt} নামে File তৈরি/খোলা হবে এবং সেখানে

\begin{center}
\textenglish{\ttfamily ICT = 85}
\end{center}

লেখা হবে। শেষে \textenglish{\ttfamily close()} File বন্ধ করবে।
\end{minipage}
\end{codeoutputbox}

%================================================
% File mode
%================================================

\booksubheading{File Mode}

File খোলার সময় Mode নির্বাচন খুব গুরুত্বপূর্ণ।

\begin{tabularx}{\linewidth}{|C{0.12\linewidth}|>{\bfseries\RaggedRight\arraybackslash}p{0.20\linewidth}|Y|}
\hline
\rowcolor{softBlue}
\textbf{Mode} & \textbf{কাজ} & \textbf{গুরুত্বপূর্ণ বিষয়}\\
\hline
\textenglish{\ttfamily r} & Read & File থেকে Data পড়া হয়; File না থাকলে Error হতে পারে।\\
\hline
\textenglish{\ttfamily w} & Write & নতুন করে Data লেখা হয়; File না থাকলে তৈরি হয়, থাকলে আগের Content মুছে নতুন Content লেখা হয়।\\
\hline
\textenglish{\ttfamily a} & Append & আগের Content ঠিক রেখে File-এর শেষে নতুন Data যোগ করা হয়; File না থাকলে তৈরি হতে পারে।\\
\hline
\textenglish{\ttfamily x} & Create & নতুন File তৈরি করে; একই নামের File আগে থেকেই থাকলে Error হয়।\\
\hline
\end{tabularx}

\begin{rememberbox}
\textbf{\color{deepOrange}MCQ Trap}

\begin{center}
\textenglish{\ttfamily r} = Read
\qquad
\textenglish{\ttfamily w} = Write/Overwrite
\qquad
\textenglish{\ttfamily a} = Append
\end{center}

বিশেষ করে \textenglish{\ttfamily w} Mode ব্যবহার করলে পুরোনো File-এর
আগের Content হারিয়ে যেতে পারে।
\end{rememberbox}

\textenglish{\ttfamily t} Text Mode এবং \textenglish{\ttfamily b} Binary
Mode বোঝাতে ব্যবহৃত হতে পারে। সাধারণ Text File-এর ক্ষেত্রে
\textenglish{\ttfamily t} হলো Default; তাই সাধারণভাবে শুধু
\textenglish{\ttfamily r}, \textenglish{\ttfamily w} বা
\textenglish{\ttfamily a} লিখলেই Text Mode বোঝায়।

%================================================
% write
%================================================

\booksubheading{File-এ Data লেখা — \textenglish{\ttfamily write()}}

\textenglish{\ttfamily write()} Method ব্যবহার করে Text File-এ String
Data লেখা যায়।

\begin{codeoutputbox}{Write Mode-এর উদাহরণ}
\noindent
\begin{minipage}[t]{0.54\linewidth}
\vspace{0pt}
\codeheading
\begin{english}
\ttfamily
file = open("student.txt", "w")\\
file.write("Rahim 85")\\
file.close()
\end{english}
\end{minipage}
\hfill
\begin{minipage}[t]{0.42\linewidth}
\vspace{0pt}
\outputheading

\textenglish{\ttfamily student.txt} File-এর Content হবে—

\begin{center}
\textenglish{\ttfamily Rahim 85}
\end{center}

\textenglish{\ttfamily w} Mode হওয়ায় File আগে থাকলে তার পুরোনো Content
নতুন Content দিয়ে প্রতিস্থাপিত হবে।
\end{minipage}
\end{codeoutputbox}

\begin{rememberbox}
\textbf{\color{deepOrange}মনে রাখবে}

\textenglish{\ttfamily write()}-এ সাধারণত String Data দিতে হয়।
Number সরাসরি লিখতে হলে প্রয়োজনে \textenglish{\ttfamily str()} দিয়ে
String-এ রূপান্তর করা যায়।
\end{rememberbox}

%================================================
% read
%================================================

\booksubheading{File থেকে Data পড়া — \textenglish{\ttfamily read()}}

\textenglish{\ttfamily r} Mode-এ File খুলে
\textenglish{\ttfamily read()} Method ব্যবহার করলে File-এর Content
পড়া যায়।

\begin{codeoutputbox}{Read Mode-এর উদাহরণ}
\noindent
\begin{minipage}[t]{0.54\linewidth}
\vspace{0pt}
\codeheading
\begin{english}
\ttfamily
file = open("student.txt", "r")\\
data = file.read()\\
print(data)\\
file.close()
\end{english}
\end{minipage}
\hfill
\begin{minipage}[t]{0.42\linewidth}
\vspace{0pt}
\outputheading

যদি File-এ থাকে—

\begin{center}
\textenglish{\ttfamily Rahim 85}
\end{center}

তবে Output হবে—

\begin{center}
\textenglish{\ttfamily Rahim 85}
\end{center}
\end{minipage}
\end{codeoutputbox}

এখানে File-এর Stored Data প্রথমে
\textenglish{\ttfamily data} Variable-এ এসেছে; এরপর
\textenglish{\ttfamily print()} দিয়ে Screen-এ দেখানো হয়েছে।

%================================================
% read methods
%================================================

\booksubheading{File পড়ার গুরুত্বপূর্ণ Method}

\begin{tabularx}{\linewidth}{|>{\bfseries\RaggedRight\arraybackslash}p{0.27\linewidth}|Y|}
\hline
\rowcolor{softBlue}
\textbf{Method} & \textbf{কাজ}\\
\hline
\textenglish{\ttfamily read()} & File-এর Content পড়ে; কোনো Size না দিলে সাধারণত অবশিষ্ট পুরো Content পড়তে পারে।\\
\hline
\textenglish{\ttfamily readline()} & একবারে একটি Line পড়তে ব্যবহৃত হয়।\\
\hline
\textenglish{\ttfamily readlines()} & File-এর Lineগুলোকে List আকারে পেতে ব্যবহার করা যায়।\\
\hline
\end{tabularx}

HSC পর্যায়ে মূল ধারণা হিসেবে
\textenglish{\ttfamily read()} এবং
\textenglish{\ttfamily readline()}-এর কাজ বুঝে রাখাই গুরুত্বপূর্ণ।

%================================================
% append
%================================================

\booksubheading{আগের Data না মুছে নতুন Data যোগ — Append Mode}

Existing File-এর Data রেখে শেষে নতুন Data যোগ করতে
\textenglish{\ttfamily a} Mode ব্যবহার করা হয়।

\begin{codeoutputbox}{Append Mode-এর উদাহরণ}
\noindent
\begin{minipage}[t]{0.54\linewidth}
\vspace{0pt}
\codeheading
\begin{english}
\ttfamily
file = open("student.txt", "a")\\
file.write(" Karim 90")\\
file.close()
\end{english}
\end{minipage}
\hfill
\begin{minipage}[t]{0.42\linewidth}
\vspace{0pt}
\outputheading

File-এ আগে যদি থাকে—

\begin{center}
\textenglish{\ttfamily Rahim 85}
\end{center}

তবে Append করার পর Content-এর শেষে নতুন Data যোগ হবে—

\begin{center}
\textenglish{\ttfamily Rahim 85 Karim 90}
\end{center}
\end{minipage}
\end{codeoutputbox}

\begin{center}
\textenglish{\ttfamily w}
\(\longrightarrow\)
\textbf{নতুন করে লেখা}
\qquad
\textenglish{\ttfamily a}
\(\longrightarrow\)
\textbf{শেষে যোগ করা}
\end{center}

%================================================
% close
%================================================

\booksubheading{\textenglish{\ttfamily close()} কেন প্রয়োজন?}

File-এর কাজ শেষ হলে \textenglish{\ttfamily close()} Method দিয়ে File
বন্ধ করা ভালো অভ্যাস। এতে File-এর সঙ্গে Program-এর ব্যবহৃত Resource
মুক্ত হয় এবং লিখিত Data সঠিকভাবে সংরক্ষিত হওয়ার প্রক্রিয়া সম্পন্ন হয়।

\begin{center}
\textbf{Open}
\(\longrightarrow\)
\textbf{Read / Write}
\(\longrightarrow\)
\textbf{Close}
\end{center}

%================================================
% with open
%================================================

\booksubheading{\textenglish{\ttfamily with open(...)} — নিরাপদ ও সহজ পদ্ধতি}

Python-এ File Handling-এর জন্য
\textenglish{\ttfamily with open(...)} ব্যবহার করলে Block-এর কাজ শেষ
হওয়ার পর File স্বয়ংক্রিয়ভাবে Close হয়ে যায়। তাই আলাদা করে
\textenglish{\ttfamily close()} লেখা সাধারণত প্রয়োজন হয় না।

\begin{codeoutputbox}{\textenglish{\ttfamily with} ব্যবহার করে File লেখা}
\noindent
\begin{minipage}[t]{0.55\linewidth}
\vspace{0pt}
\codeheading
\begin{english}
\ttfamily
with open("result.txt", "w") as file:\\
\hspace*{4mm}file.write("Average = 85")
\end{english}
\end{minipage}
\hfill
\begin{minipage}[t]{0.41\linewidth}
\vspace{0pt}
\outputheading

\textenglish{\ttfamily result.txt}-এ লেখা হবে—

\begin{center}
\textenglish{\ttfamily Average = 85}
\end{center}

Block শেষ হলে File স্বয়ংক্রিয়ভাবে Close হবে।
\end{minipage}
\end{codeoutputbox}

\begin{rememberbox}
\textbf{\color{deepOrange}Shortcut}

\begin{center}
\textenglish{\ttfamily open() + close()}
\(\longrightarrow\)
\textbf{Manual File Management}

\medskip

\textenglish{\ttfamily with open(...)}
\(\longrightarrow\)
\textbf{Automatic Close}
\end{center}
\end{rememberbox}

%================================================
% Complete save and read
%================================================

\booksubheading{একই Program-এ Data Save ও পুনরায় Read}

নিচের উদাহরণে প্রথমে Result File-এ Save করা হয়েছে। পরে একই File আবার
Open করে Stored Data পড়া হয়েছে। এতে Python Program ও Persistent File
Data-এর সম্পর্ক স্পষ্ট হয়।

\begin{codeoutputbox}{Result Save করে পরে Read}
\noindent
\begin{minipage}[t]{0.58\linewidth}
\vspace{0pt}
\codeheading
\begin{english}
\ttfamily
marks = 85\\[0.4mm]
with open("result.txt", "w") as file:\\
\hspace*{4mm}file.write(str(marks))\\[0.6mm]
with open("result.txt", "r") as file:\\
\hspace*{4mm}saved = file.read()\\[0.6mm]
print(saved)
\end{english}
\end{minipage}
\hfill
\begin{minipage}[t]{0.38\linewidth}
\vspace{0pt}
\outputheading

\begin{center}
\textenglish{\ttfamily 85}
\end{center}

ধাপগুলো—

\begin{enumerate}
\item \textenglish{\ttfamily marks}-এ 85 রাখা হয়েছে।
\item \textenglish{\ttfamily result.txt}-এ 85 Save হয়েছে।
\item File আবার Read Mode-এ খোলা হয়েছে।
\item Saved Data পড়ে Output দেখানো হয়েছে।
\end{enumerate}
\end{minipage}
\end{codeoutputbox}

\begin{center}
\textbf{Variable}
\(\longrightarrow\)
\textbf{\textenglish{\ttfamily write()}}
\(\longrightarrow\)
\textbf{File-এ Persistent Data}
\(\longrightarrow\)
\textbf{\textenglish{\ttfamily read()}}
\(\longrightarrow\)
\textbf{আবার Program-এ Data}
\end{center}

%================================================
% File handling with function
%================================================

\booksubheading{Function ও File একসঙ্গে ব্যবহার}

Function Data Processing করতে পারে এবং Process করা Result File-এ
সংরক্ষণ করা যায়। অর্থাৎ Function ও File প্রতিদ্বন্দ্বী নয়; Program-এ
দুটির কাজ আলাদা হলেও একসঙ্গে ব্যবহার করা যায়।

\begin{codeoutputbox}{Function-এর Result File-এ Save}
\noindent
\begin{minipage}[t]{0.58\linewidth}
\vspace{0pt}
\codeheading
\begin{english}
\ttfamily
def average(a, b):\\
\hspace*{4mm}return (a + b) / 2\\[0.5mm]
result = average(80, 90)\\[0.5mm]
with open("average.txt", "w") as file:\\
\hspace*{4mm}file.write(str(result))
\end{english}
\end{minipage}
\hfill
\begin{minipage}[t]{0.38\linewidth}
\vspace{0pt}
\outputheading

Function Result—

\begin{center}
\textenglish{\ttfamily 85.0}
\end{center}

\textenglish{\ttfamily average.txt} File-এও
\textenglish{\ttfamily 85.0} সংরক্ষিত হবে।

\medskip
Function \(\rightarrow\) Processing

File \(\rightarrow\) Persistent Storage
\end{minipage}
\end{codeoutputbox}

%================================================
% File handling mistakes
%================================================

\booksubheading{Python File Handling-এর সাধারণ ভুল}

\begin{rememberbox}
\begin{itemize}
\item \textenglish{\ttfamily w} এবং \textenglish{\ttfamily a} Mode-কে একই মনে করা।
\item File না থাকা সত্ত্বেও \textenglish{\ttfamily r} Mode-এ Open করলে কী হতে পারে তা না বোঝা।
\item \textenglish{\ttfamily open()} করার পর File Object-এর বদলে শুধু File Name নিয়েই কাজ করার চেষ্টা করা।
\item কাজ শেষে \textenglish{\ttfamily close()} না করা, যখন \textenglish{\ttfamily with} ব্যবহার করা হয়নি।
\item Number সরাসরি \textenglish{\ttfamily write()}-এ দিলে সমস্যা হতে পারে—প্রয়োজনে \textenglish{\ttfamily str()} ব্যবহার করতে হয়।
\item \textenglish{\ttfamily .py} Program File এবং \textenglish{\ttfamily .txt} Data File-কে একই মনে করা।
\end{itemize}
\end{rememberbox}

%================================================
% File handling quick revision
%================================================

\booksubheading{File Handling এক নজরে}

\begin{tabularx}{\linewidth}{|>{\bfseries\RaggedRight\arraybackslash}p{0.31\linewidth}|Y|}
\hline
\rowcolor{softBlue}
\textbf{বিষয়} & \textbf{মনে রাখবে}\\
\hline
File Open & \textenglish{\ttfamily open("name", "mode")}\\
\hline
Read Mode & \textenglish{\ttfamily r}\\
\hline
Write Mode & \textenglish{\ttfamily w}\\
\hline
Append Mode & \textenglish{\ttfamily a}\\
\hline
Data Read & \textenglish{\ttfamily read()}, \textenglish{\ttfamily readline()}\\
\hline
Data Write & \textenglish{\ttfamily write()}\\
\hline
File Close & \textenglish{\ttfamily close()}\\
\hline
Automatic Close & \textenglish{\ttfamily with open(...)}\\
\hline
Persistent Data & Program বন্ধ হওয়ার পরও File-এ সংরক্ষিত Data\\
\hline
\end{tabularx}
%================================================
% Practical flow
%================================================

%================================================
% Practical flow
%================================================

\booksubheading{বাস্তব উদাহরণ}

Student Result Program-এ—

\begin{center}

\textbf{Marks Input}
\(\longrightarrow\)
\textbf{Result Calculate}
\(\longrightarrow\)
\textbf{Variable-এ Result}

\end{center}

শুধু Variable ব্যবহার করলে Program বন্ধ হওয়ার পর Runtime Data থাকবে না।

কিন্তু—

\begin{center}

\textbf{Marks Input}
\(\longrightarrow\)
\textbf{Result Calculate}
\(\longrightarrow\)
\textbf{File-এ Save}
\(\longrightarrow\)
\textbf{পরে আবার ব্যবহার}

\end{center}

করলে Data ভবিষ্যতেও ব্যবহার করা সম্ভব।

%================================================
% Function vs File
%================================================

\booksubheading{Function ও File-এর কাজের পার্থক্য}

\begin{tabularx}
{\linewidth}
{
|>{\bfseries\RaggedRight\arraybackslash}p{0.25\linewidth}
|Y
|Y|
}

\hline
\rowcolor{softBlue}

\textbf{বিষয়}
&
\textbf{Function}
&
\textbf{File}
\\
\hline

প্রধান কাজ
&
Program-এর Logic সংগঠিত করে
&
Data সংরক্ষণ করে
\\
\hline

Input
&
Data গ্রহণ করতে পারে
&
Stored Data ধারণ করে
\\
\hline

Processing
&
Calculation/Logic করতে পারে
&
মূল উদ্দেশ্য Storage
\\
\hline

ফলাফল
&
Value Return করতে পারে
&
Data পরে পুনরায় ব্যবহার করা যায়
\\
\hline

\end{tabularx}

\begin{center}

\textbf{Function}
\(\longrightarrow\)
\textbf{কাজ/Logic সংগঠিত করে}

\qquad

\textbf{File}
\(\longrightarrow\)
\textbf{Data সংরক্ষণ করে}

\end{center}

%================================================
% Combined Example
%================================================

\booksubheading{একই উদাহরণে পুরো ধারণা}

\begin{codeoutputbox}{Average নির্ণয়ের Function}

\noindent
\begin{minipage}[t]{0.55\linewidth}
\codeheading

\begin{english}
\ttfamily
def average(a, b):\\
\hspace*{4mm}result = (a + b) / 2\\
\hspace*{4mm}return result\\[0.5mm]
x = average(80, 90)\\
print(x)
\end{english}

\end{minipage}
\hfill
\begin{minipage}[t]{0.41\linewidth}
\outputheading

\begin{center}
\textenglish{\ttfamily 85.0}
\end{center}

\textenglish{\ttfamily a, b}
\(\longrightarrow\)
Input / Parameter

\medskip

\textenglish{\ttfamily (a+b)/2}
\(\longrightarrow\)
Processing

\medskip

\textenglish{\ttfamily result}
\(\longrightarrow\)
Local Variable

\medskip

\textenglish{\ttfamily 85.0}
\(\longrightarrow\)
Return Value

\end{minipage}

\end{codeoutputbox}

\begin{tabularx}
{\linewidth}
{
|>{\bfseries\RaggedRight\arraybackslash}p{0.34\linewidth}
|Y|
}

\hline
\rowcolor{softBlue}

\textbf{বিষয়}
&
\textbf{উদাহরণ}
\\
\hline

Function
&
\textenglish{\ttfamily average()}
\\
\hline

Contract-এর Input
&
\textenglish{\ttfamily a, b}
\\
\hline

Processing
&
\textenglish{\ttfamily (a + b) / 2}
\\
\hline

Local Variable
&
\textenglish{\ttfamily result}
\\
\hline

Argument
&
\textenglish{\ttfamily 80, 90}
\\
\hline

Return Value
&
\textenglish{\ttfamily 85.0}
\\
\hline

\end{tabularx}

যদি
\textenglish{\ttfamily 85.0}
শুধু Program চলাকালে প্রয়োজন হয়, Variable যথেষ্ট।

কিন্তু Program বন্ধ হওয়ার পরও Result রাখতে হলে
\textbf{File}-এ সংরক্ষণ করতে হবে।

%================================================
% MCQ Points
%================================================

\booksubheading{গুরুত্বপূর্ণ MCQ তথ্য}

\begin{rememberbox}

\begin{tabularx}
{\linewidth}
{
>{\bfseries\RaggedRight\arraybackslash}p{0.38\linewidth}
Y
}

Computational Thinking
&
সমস্যাকে যৌক্তিকভাবে বিশ্লেষণ ও সমাধান পরিকল্পনা
\\[0.45mm]

Function Contract
&
Input + Processing + Expected Output
\\[0.45mm]

Scope
&
Variable কোথায় ব্যবহার করা যাবে
\\[0.45mm]

Local Variable
&
Function-এর ভিতরে তৈরি
\\[0.45mm]

Local Scope
&
Local Variable-এর ব্যবহারযোগ্য সীমা
\\[0.45mm]

Global Variable
&
Function-এর বাইরে তৈরি
\\[0.45mm]

Global Scope
&
Global Variable-এর বিস্তৃত ব্যবহারযোগ্য সীমা
\\[0.45mm]

Parameter
&
Function-এর Local Scope-এর অংশ
\\[0.45mm]

Global Value Read
&
সাধারণত
\textenglish{\ttfamily global}
Keyword প্রয়োজন নেই
\\[0.45mm]

Global Assignment
&
\textenglish{\ttfamily global}
Keyword প্রয়োজন হতে পারে
\\[0.45mm]

Temporary Data
&
Runtime চলাকালীন সাময়িক Data
\\[0.45mm]

Persistent Data
&
Program শেষ হওয়ার পরও সংরক্ষিত Data
\\[0.45mm]

File
&
Persistent Data সংরক্ষণে ব্যবহৃত হয়
\\[0.45mm]

File Handling
&
File Open, Read/Write/Append ও Close করার প্রক্রিয়া
\\[0.45mm]

\textenglish{\ttfamily open()}
&
File খুলে File Object তৈরি করে
\\[0.45mm]

\textenglish{\ttfamily r / w / a}
&
Read / Write (Overwrite) / Append Mode
\\[0.45mm]

\textenglish{\ttfamily read() / write()}
&
File থেকে Data পড়া / File-এ Data লেখা
\\[0.45mm]

\textenglish{\ttfamily with open(...)}
&
Block শেষে File স্বয়ংক্রিয়ভাবে Close হয়
\\

\end{tabularx}

\end{rememberbox}

%================================================
% Common Mistakes
%================================================

\booksubheading{সাধারণ ভুল}

\begin{rememberbox}

\begin{itemize}

\item Function Contract-কে শুধু Function-এর নাম মনে করা।

\item Scope-কে Variable-এর Value মনে করা।

\item Local Variable Function-এর বাইরে ব্যবহার করা যাবে মনে করা।

\item Local ও Global Variable একই মনে করা।

\item Function-এর ভিতরে একই Name-এ Assignment করলেই Global Variable
পরিবর্তিত হবে মনে করা।

\item Global Variable শুধু পড়ার জন্যও সবসময়
\textenglish{\ttfamily global}
লাগে মনে করা।

\item
\textenglish{\ttfamily global}
Keyword-এর কাজ ভুল বোঝা।

\item সব Data Global Variable-এ রাখা ভালো Design মনে করা।

\item Variable-এর Data Program বন্ধ হওয়ার পরও স্থায়ী থাকবে মনে করা।

\item Temporary Data ও Persistent Data-কে একই মনে করা।

\item Runtime Memory ও File-কে একই ধরনের Storage মনে করা।

\end{itemize}

\end{rememberbox}

%================================================
% Final Reminder
%================================================

\begin{rememberbox}

\textbf{\color{deepOrange}মনে রাখবে}

\begin{tabularx}
{\linewidth}
{
>{\bfseries\RaggedRight\arraybackslash}p{0.25\linewidth}
Y
>{\bfseries\RaggedRight\arraybackslash}p{0.25\linewidth}
Y
}

Function Contract
&
Input + Processing + Output
&
Scope
&
Variable-এর Access সীমা
\\[0.5mm]

Local Variable
&
Function-এর ভিতরে
&
Global Variable
&
Function-এর বাইরে
\\[0.5mm]

Local Scope
&
Local-এর ব্যবহারযোগ্য সীমা
&
Global Scope
&
Global-এর ব্যবহারযোগ্য সীমা
\\[0.5mm]

Parameter
&
Local Scope
&
\textenglish{\ttfamily global}
&
Global Variable-এ Assignment
\\[0.5mm]

Parameter + Return
&
পরিষ্কার Function Design
&
Temporary Data
&
Runtime Data
\\[0.5mm]

File
&
Persistent Storage
&
Persistent Data
&
পরে আবার ব্যবহারযোগ্য
\\[0.5mm]

\textenglish{\ttfamily open()}
&
File Open / File Object
&
\textenglish{\ttfamily r, w, a}
&
Read, Write, Append
\\[0.5mm]

\textenglish{\ttfamily read()} / \textenglish{\ttfamily write()}
&
Data পড়া / লেখা
&
\textenglish{\ttfamily with open(...)}
&
Automatic Close
\\

\end{tabularx}

\medskip

\begin{center}

\textbf{Function কী নেবে/দেবে?}
\(\longrightarrow\)
\textbf{Contract}

\qquad

\textbf{Variable কোথায় ব্যবহার হবে?}
\(\longrightarrow\)
\textbf{Scope}

\medskip

\textbf{Function-এর নিজের Variable}
\(\longrightarrow\)
\textbf{Local}

\qquad

\textbf{Function-এর বাইরের Variable}
\(\longrightarrow\)
\textbf{Global}

\medskip

\textbf{Program চলাকালীন Data}
\(\longrightarrow\)
\textbf{Temporary}

\qquad

\textbf{Program বন্ধের পরও Data}
\(\longrightarrow\)
\textbf{Persistent / File}

\medskip

\textbf{Python দিয়ে File ব্যবহার}
\(\longrightarrow\)
\textbf{\textenglish{\ttfamily open()} \(\rightarrow\) Read/Write/Append \(\rightarrow\) Close}

\end{center}

\end{rememberbox}

\small

%================================================
% 5.46
%================================================

